\documentclass{amsart}
\usepackage{lmodern}
\usepackage[english]{babel}
\usepackage{hyperref}
\usepackage{cite}
\usepackage{cleveref}
\usepackage[all]{xy}
\usepackage{hhline}
\usepackage{multirow}
\usepackage{mathtools}
\usepackage{listings}

\lstset{
  basicstyle=\ttfamily,
  mathescape
}

\theoremstyle{definition}
\newtheorem{ex}{Example}

\renewcommand{\theenumi}{\alph{enumi}}

\def\CC{\mathbb{C}}
\def\KK{\mathbb{K}}
\def\PP{\mathbb{P}}
\def\NN{\mathbb{N}}

\def\dA{{\ttfamily DiffAlg}}

%%%%%%%%%%%%%%%%%%%%%%%%%%%%%%%%%%%%%%%%%%%%%%%%%%%%%%%%%%%%%%%%%%%%%%%%%%
%%%%%%%%%%%%%%%%%%%%%%%%%%%%%%%%%%%%%%%%%%%%%%%%%%%%%%%%%%%%%%%%%%%%%%%%%%
%%%%%%%%%%%%%%%%%%%%%%%%%%%%%%%%%%%%%%%%%%%%%%%%%%%%%%%%%%%%%%%%%%%%%%%%%%
%%%%%%%%%%%%%%%%%%%%%%%%%%%%%%%%%%%%%%%%%%%%%%%%%%%%%%%%%%%%%%%%%%%%%%%%%%
%%%%%%%%%%%%%%%%%%%%%%%%%%%%%%%%%%%%%%%%%%%%%%%%%%%%%%%%%%%%%%%%%%%%%%%%%%
%%%%%%%%%%%%%%%%%%%%%%%%%%%%%%%%%%%%%%%%%%%%%%%%%%%%%%%%%%%%%%%%%%%%%%%%%%

\title{\dA: a Differential Algebra Package}

\author[M. Dubinsky]{Manuel Dubinsky}
\address{Departamento de Computaci\'on\\ Universidad de Avellaneda\\ Argentina}
\email{manudubinsky@gmail.com}

\author[C. Massri]{C\'esar Massri}
\address{Departamento de Matem\'atica\\ Universidad CAECE\\ Argentina}
\email{cmassri@caece.edu.ar}

\author[A. Molinuevo]{Ariel Molinuevo}
\address{Instituto de Matem\'atica\\
Universidade Federal do Rio de Janeiro\\
Brasil}
\email{amoli@dm.uba.ar}


\author[F. Quallbrunn]{Federico Quallbrunn}
\address{Departamento de Matem\'atica\\ Universidad CAECE y FCEN, UBA\\ Argentina}
\email{fquallb@dm.uba.ar}

\thanks{The authors were supported by CONICET, Argentina.}


\subjclass[2010]{53-04; 14-04}
\keywords{Differential operators, Exterior algebra, Vector Fields}


\begin{document}

\begin{abstract}
In this article we present \dA, a differential algebra package for Macaulay2. 
It can perform the following operations: wedge products and exterior differential of differential forms, contraction and Lie derivative of a differential form with respect to a vector field and Lie brackets between vector fields.

Given a homogeneous differential operator of degree one $D$, the lack of an algebraic module structure attached to the kernel or image of $D$ hinders the study of $D$. The main purpose of \dA\ is to handle these spaces degree-wise.
\end{abstract}

\maketitle

\section{Motivation and description of the package}

\

Algebraic and differential operations arise naturally when working with differential forms and vector fields, \emph{e.g.}:  wedge products and exterior differential of differential forms, contraction and Lie derivative of a differential form with respect to a vector field and Lie brackets between vector fields. Some important statements involving these operations include the following:
\begin{enumerate}
\item A differential $r$-form $\omega$ in the affine space $\KK^{n+1}$ \emph{descends to the projective space} $\PP^n_{\KK}$, if it satisfies the equation
\[
i_R\omega=0,
\]
where $\KK$ is a field, $R$ is the \text{radial vector field} $R=\sum x_i\frac{\partial}{\partial x_i}$ and $i_R$ denotes the contraction, see \cite[Theorem 8.13, p.~176]{hart}.
\item If $\omega$ is a differential $1$-form, then $\omega$ defines a foliation in $\KK^{n+1}$ if it satisfies the \emph{Frobenius integrability condition} given by the equation
\[
\omega\wedge d\omega=0,
\]
see \cite[Definition 2.2, p.~823]{suwa-review}.
\item Let $\omega$ be an integrable 1-form and $L_X\omega$ the Lie derivative of $\omega$ with respect to a vector field $X$. Then, the solutions of the equation 
\[
L_X\omega=0
\]
define all the \emph{infinitesimal automorphisms} of the foliation given by $\omega$,
see \cite[Proposition 7.7, p.~845]{suwa-review}.
\item Let $\omega$ be an integrable 1-form. Then the \emph{tangent space of the space of foliations} at $\omega$ is given by the differential $1$-forms $\eta$ that satisfy the equation
\[
\omega\wedge d\eta + d\omega\wedge \eta=0,
\]
see \cite[Section 2.1. p.~709]{fji}.
\item Let $D$ be a \emph{bracket generating distribution}.
Some important invariants of $D$ are the ranks of the derived sequence
\[
 a(p):=\text{rank}\, D^{(p)}=\mathrm{rank}\left( D^{(p-1)}+[D,D^{(p-1)}]\right),
 \]
 see \cite[\S 1, pp.~8-9]{Tanaka}.
\item A \emph{symplectic structure} in a variety of dimension $2r$ is given by a $2$-form $\omega$ such that $d\omega=0$ and $\omega^{r}\neq 0$, see \cite[p.~41]{bcggg}.
% \item The exterior algebra of differential forms of the projective space $\PP^n_\KK$ can be endowed with a graded, associative and super commutative product given by the \emph{second multiplication of Gelfand-Kapranov-Zelevinsky} defined as
% \[
% \omega\star\eta =
% \left\{
% \begin{aligned}
% &\frac{e}{e+d}\ \omega\wedge d\eta - (-1)^r \frac{d}{e+d}\ d\omega\wedge\eta\qquad &&\text{if }e,d\neq 0 \\
% &0&&\text{if $e$ or $d$ equals $0$}
% \end{aligned}
% \right.
% \]
% where $\omega$ is a differential $r$-form and $\eta$ is a differential $s$-form with homogeneous coefficients of degree $e-r$ and $d-s$ respectively, see \cite[Chapter 2, 4. B., p.~72]{gkz}.
\end{enumerate}

\bigskip

For a clear understanding on how \dA\ deals with such equations, let us fix some notation.

\bigskip

Let $S=\KK[x_0,\ldots,x_n]$ be the polynomial ring in $n+1$ variables and let $\Omega=\bigoplus_{r\geq 0} \Omega^r$ be the exterior algebra of differential forms of $S$ over $\KK$. Let $\Omega^r(d)$ denote the space of $r$-forms with polynomial coefficients of homogeneous degree $d$, where we assign degree 1 to each $x_i$ and degree 0 to each $dx_i$.

Therefore $\omega\in\Omega^r(d)$ can be written as
\begin{equation}\label{omega}
\omega = \sum_{\substack{I\subset \{0,\ldots,n\} \\ \#I=r}} \sum_{\substack{\alpha\in\NN^{n+1}\\ |\alpha|=d}} \ a_{\alpha,I}\ x^\alpha \ dx_I,\qquad a_{\alpha,I}\in\KK,
\end{equation}
where for each $I$ of the form $I=\{i_1,\ldots,i_r\}\subset \{0,\ldots,n\}$ let $dx_I$ denote $dx_{i_1}\wedge\ldots\wedge dx_{i_r}$ and each $\alpha=(\alpha_0,\ldots,\alpha_n)\in\NN^{n+1}$ denote $|\alpha|:=\sum_{i=0}^n\alpha_i$ and $x^\alpha:=x_0^{\alpha_0}.\ldots.x_n^{\alpha_n}$.

\bigskip

Let $T$ be the module of vector fields with coefficients in $S$. Let $T(e)$ denote the homogeneous vector fields with polynomial coefficients of degree $e$, where we assign degree 1 to each $x_i$ and degree 0 to each $\frac{\partial}{\partial x_i}$ and, analogously to \cref{omega}, $X\in T(e)$ can be written as
\[
X = \sum_{i=0}^n \sum_{\substack{\beta\in\NN^{n+1}\\ |\beta|=e}} \ b_{\beta,i}\ x^\beta \ \frac{\partial}{\partial x_i},\qquad b_{\beta,i}\in\KK.
\]

\bigskip

Current algebraic software systems implement functionality to deal with differential forms and vector fields, but usually scalar parameters $a_{\alpha,I}$ and $b_{\beta,i}$ must be specified as fixed elements in $\KK$.
Instead, \dA\ treats homogeneous forms and vector fields in a completely symbolic environment by considering the scalar coefficient rings
\[
\KK[a_{\alpha,I}]\qquad \text{and}\qquad \KK[b_{\beta,i}].
\]
Scalar coefficients can be systematically obtained by looking at the coordinates of differential forms and vector fields written in the standard basis
\[
\mathcal{B}_{r,d} = \left\{x^\alpha dx_I\right\}_{\substack{\#I=r\\|\alpha|=d}}\qquad\text{and}\qquad \mathcal{B}_{e} = \left\{ x^\beta\tfrac{\partial}{\partial x_i}\right\}_{|\beta|=e}
\]
of the spaces $\Omega^r(d)$ and $T(e)$, respectively.

\bigskip


Importantly, when using \dA, each object is expected to be defined in its own coefficient ring. Then, computing some operations like wedge product or contraction
involve different input and output rings, producing a modification of the coefficients rings $\KK[a_{\alpha,I}]$ or $\KK[b_{\beta,i}]$. For greater clarity, consider the following example. Fix $\omega\in\Omega^r(d)$ and $X\in T(e)$ and consider the contraction $i_X\omega\in\Omega^{r-1}(d+e)$.
Then, the following elements will be taking place:

\[
\begin{tabular}{ l  c  c c c }
  \small{Operation} & $(\omega,X)$ & $\xmapsto{\hspace{1cm}}$& $i_X\omega$\\[.9em]
  \small{Rings} & $\KK[a_{\alpha,I}]\times\KK[b_{\beta,i}]$ & & $\KK[a_{\alpha,I},b_{\beta,i}]$\\[.6em]
  \small{Basis} & $\mathcal{B}_{r,d}\times\mathcal{B}_{e}$ & &  $\mathcal{B}_{r-1,d+e}$ \\[.6em]
\end{tabular}
\]

\bigskip

As mentioned before, the main purpose of \dA\ is to find algebraic solutions to equations in the context of differential algebra. Equations are treated differently in the linear and non-linear cases:
\begin{enumerate}
\item[i)] In the linear case, for example $i_R\omega=0$, \dA\ can compute a basis of the solutions of the equation. Once this is done, it can also compute a generic linear combination of the elements of the basis, see Example \ref{ex1}.
\item[ii)] In the non-linear case, for example $\omega\wedge d\omega=0$, the coordinates will be polynomial. In this case \dA\ would compute the ideal generating the space of solutions. This ideal can be obtained in two different ways: taking coordinates in the basis $\mathcal{B}_{r,d}$ or $\mathcal{B}_{e}$ or taking coordinates in the basis $\{dx_I\}$ or $\{\frac{\partial}{\partial x_i}\}$, see Examples \ref{ex2} and \ref{ex4}.
\end{enumerate}

\bigskip
\dA\ can also be a valuable tool for studying differential operators. The lack of an algebraic theory to deal with such objects can be mitigated by non-conclusive computations easily made by \dA. As a first example, one could consider computing solutions of the annihilation of a differential operator degree-wise for low degrees.


\section{Some examples}

\ex\label{ex1} In the following example we obtain a basis of the space of projective differential 2-forms in $\PP_\KK^3$. Then, we define a generic projective differential 2-form to be possibly used in further computations.

\bigskip

{\tt\footnotesize
\begin{lstlisting}
i1 : loadPackage "DiffAlg";

i2 : R = radial 3

o2 = x ax  + x ax  + x ax  + x ax
      0  0    1  1    2  2    3  3

o2 : DiffAlgField

i3 : w = newForm(3,2,1,"a");   

o3 = (a x +a x +a  x +a  x )dx dx +(a x  + a x  + a  x  + a  x )dx dx 
     0 0  6 1  12 2  18 3   0  1   1 0    7 1    13 2    19 3   0  2
     ---------------------------------------------------------------
     + (a x  + a x  + a  x  + a  x )dx dx  + (a x  + a x  + a  x  +
         3 0    9 1    15 2    21 3   1  2     2 0    8 1    14 2  
     ---------------------------------------------------------------
     a  x )dx dx  + (a x  + a  x  + a  x  + a  x )dx dx +(a x +a  x  +
      20 3   0  3     4 0    10 1    16 2    22 3   1  3   5 0  11 1  
     ---------------------------------------------------------------
     a  x  + a  x )dx dx
      17 2    23 3   2  3

o3 : DiffAlgForm

i4 : pretty ring w

      QQ[i]
o4 = ------[][a , a , a , a , a , a , a , a , a , a , a  , a  , a  ,
      2        0   1   2   3   4   5   6   7   8   9   10   11   12 
     i  + 1
     ---------------------------------------------------------------
      a  , a  , a  , a  , a  , a  , a  , a  , a  , a  , a  ][x , x , 
       13   14   15   16   17   18   19   20   21   22   23   0   1 
     ---------------------------------------------------------------
     x , x ][dx , dx , dx , dx ]
      2   3    0    1    2    3
                                                                
i5 : K = genKer (R _ w, w);

i6 : length K

o6 = 4

i7 : v = linearComb(K,"a")

o7 = (a x  - a x )dx dx +(- a x  + a x )dx dx  + (a x  + a x )dx dx +
       0 2    1 3   0  1     0 1    2 3   0  2     0 0    3 3   1  2
     ----------------------------------------------------------------
     (a x  - a x )dx dx  + (-a x  - a x )dx dx  + (a x  + a x )dx dx
       1 1    2 2   0  3      1 0    3 2   1  3     2 0    3 1   2  3

o7 : DiffAlgForm

i8 : pretty ring v

      QQ[i]
o8 = ------[][a , a , a , a ][x , x , x , x ][dx , dx , dx , dx ]
       2        0   1   2   3   0   1   2   3    0    1    2    3
      i  + 1
\end{lstlisting}
}

\noindent Let us explain part of the code:
\begin{itemize}
\item[\tt i2.]
Creates the radial vector field in 4 variables. We are 
denoting the basic field ${\partial}/{\partial x_i}$ as $ax\_i$.
\item[\tt i3.]
Creates a generic linear 2-form in $\Omega^2_{\KK^4}(1)$, where the coefficients 
are indexed as $a_i$.
\item[\tt i4.]
Shows the ring of definition of $w$.
\item[\tt i5.]
Gets a basis (as a Macaulay2 list) of forms in $\Omega^2_{\KK^4}(1)$ that 
descend to projective space. The operation $R \_ w$ computes the 
contraction of the differential form $w$ with the vector field $R$.
\item[\tt i6.]
Gets the dimension of $\Omega^2_{\PP^3}(1)$ in projective 3-space
\item[\tt i7.]
Defines a generic projective form with coefficients $a_i$.
\item[\tt i8.]
Shows the ring of definition of $v$.
\end{itemize}


\bigskip

\ex\label{ex2} In the finite dimensional $\KK$-vector space $\Omega^1(d)$, the solutions of the equation $\omega\wedge d\omega=0$ determine an algebraic variety, its points are the integrable differential 1-forms of degree $d$.  In the following example, we compute the equations of the variety of integrable 1-forms of degree 1 in 3-dimensional space.

It is worth mentioning that, for $n\geq 3$ and $d>5$, it is an open problem to classify the irreducible components of this varieties, see \cite{fji}.

\bigskip

{\tt\footnotesize
\begin{lstlisting}
i1 : loadPackage "DiffAlg";

i2 : w = newForm (2,1,1,"a")

o2 = (a x  + a x  + a x )dx  + (a x  + a x  + a x )dx +(a x  + a x +
       0 0    3 1    6 2   0     1 0    4 1    7 2   1   2 0    5 1    
     ---------------------------------------------------------------
        a x )dx
         8 2   2

o2 : DiffAlgForm

i3 : moduliIdeal (w ^ (diff w))

o3 = ideal (- a a  + a a  + a a  - a a , - a a  + a a  + a a  - a a ,
               2 3    0 5    1 6    0 7     2 4    1 5    4 6    3 7
     ----------------------------------------------------------------
            a a  - a a  + a a  - a a )
             5 6    2 7    1 8    3 8

               QQ[i]
o3 : Ideal of ------[][a , a , a , a , a , a , a , a , a ]
               2        0   1   2   3   4   5   6   7   8
              i  + 1
\end{lstlisting}
}

\noindent About the code:
\begin{itemize}
\item[\tt i2.]
Creates a generic linear 1-form in $\Omega^1_{\KK^3}(1)$, with coefficients $a_i$.
\item[\tt i3.]
Returns the ideal of the scalar coefficients given by 
$w \wedge d w = 0$.
\end{itemize}

\ex\label{ex3} Let $D$ be a 2-dimensional distribution generated by vector fields $X$ and $Y$ in 5-dimensional space.
In the following example we compute the ranks of the derived distributions $D^{(p)}$.
We verify that this derived series
eventually spans the entire tangent space. A distribution $D$ satisfying this
condition is called \emph{bracket-generating}.

\bigskip

{\tt\footnotesize
\begin{lstlisting}
i2 : X = newField "x_0^2*ax_0+x_1^2*ax_1+x_2^2*ax_2+x_3^2*ax_3";

i3 : Y = newField "x_5*ax_0+x_4*ax_1+x_3*ax_2+x_2*ax_3+x_1*ax_4+
         x_0*ax_5";

i4 : D_0 = {X,Y};

i5 : for b in 1..3 do (
      for a in D_(b-1) do (
       D_b = join(D_(b-1),{a|Y,a|X})
      )
     );

i6 : {rank dist D_0, rank dist D_1, rank dist D_2, rank dist D_3}

o6 = {2, 3, 5, 6}

o6 : List
\end{lstlisting}
}

\noindent About the code:
\begin{itemize}
\item[\tt i5.]
Computes the derived sequence.
\item[\tt i6.]
Prints the ranks of the derived series.
\end{itemize}

\bigskip

\ex\label{ex4} In the following example, we generate a random rational 1-form of type (1,2) in $\PP^2_\KK$. First, we compute (the dimension of) the space of its integrating factors, see \cite[pp.~828-829]{suwa-review}. Then, we compute the ideal of its singular 
locus (the ideal where it vanishes).

\bigskip

{\tt\footnotesize
\begin{lstlisting}
i2 : w = random logarithmicForm (2,{1,2},"a",Projective => true);

i3 : f = newForm (2,0,3,"a");

i4 : length genKer(w^(diff f) + f*(diff w), f)

o4 = 2

i5 : I = singularIdeal w

                         2                        2    2
o6 = ideal (- 9x x  + 63x  - 36x x  - 54x x  - 54x , 9x  - 63x x  -
                0 1      1      0 2      1 2      2    0      0 1
     -----------------------------------------------------------------
                  2     2               2
      27x x  -45x x  - 54x , 36x  + 81x x  + 45x  + 54x x  + 54x x )
         0 2     1 2      2     0      0 1      1      0 2      1 2

               QQ[i]
o6 : Ideal of ------[][x , x , x ]
               2        0   1   2
              i  + 1
\end{lstlisting}
}

% % % % % % % % % % % %
% % % % % % % % % % % %
% BIBLIOGRAPHY
% % % % % % % % % % % %
% % % % % % % % % % % %

\bibliography{biblio}{}
\bibliographystyle{amsplain}

% \bigskip

% \bigskip

\end{document}
